\customchapter{RESUMEN}
\noindent
El patrimonio pétreo está expuesto al biodeterioro microbiano, proceso mediante el cual las comunidades que colonizan la roca disuelven su matriz mineral y alteran su superficie. La Estela de Raimondi, litoescultura de granito de la cultura Chavín bajo custodia museística, presenta biopelículas y alteraciones cromáticas cuya evaluación previa se limitó a la identificación taxonómica, sin esclarecer la capacidad de deterioro de los organismos vivos que la habitan. El objetivo de este proyecto fue diseñar un flujo de análisis computacional para caracterizar el potencial de biodeterioro y las señales de adaptación genómica de los linajes bacterianos dominantes de la fracción viable del monumento, con el fin de sustentar la evaluación del riesgo y una conservación preventiva basada en evidencia genómica. Se partió de datos de secuenciación metagenómica masiva obtenidos de la fracción cultivable de la superficie, sobre los que se estructuró un flujo reproducible que integró el control de calidad genómica, la clasificación taxonómica basada en el genoma completo, la agrupación de genomas redundantes en linajes, la reconstrucción pangenómica comparativa frente a referencias globales y la anotación funcional del potencial metabólico.\\


\noindent \textbf{Palabras clave:}\\
\noindent Biodeterioro; Patrimonio pétreo; Metagenómica; Pangenoma;
Potencial metabólico; Conservación preventiva; Estela de Raimondi